% !TeX root = ./rapportUTT.tex
% !BIB TS-program = biber

\documentclass{rUTT}
% Pour retirer le thème couleur UTT,
%   Commenter la ligne précédente
%   Décommenter la ligne dessous
% \documentclass[noUTTcolors]{rUTT}


%%%%%%%%%%%%%%%%%%%%%%%%%%%%%%%%%%%%%%%%%%%%%%%%%%%%%%%%%%%%
% Quelques trucs à savoir pour modifier en paix :
% TOC = Table of Contents
% LOF / LOT = List Of Figures / List Of Tables
%%%%%%%%%%%%%%%%%%%%%%%%%%%%%%%%%%%%%%%%%%%%%%%%%%%%%%%%%%%%git@github.com:n3rada/ScribUTT.git
% Si on est on mode rapport de stage :
% Pour la confidentialité

\RPeda{{\sc PLOIX} Alain} % Responsable pédagogique

\Entreprise{Winamax} % Nom de l'entreprise / organisme / institution
\Lieu{136bis rue Grenelle, 75007 Paris} 
\REntre{{\sc FIDEL} Nicolas} % Responsable entreprise

% Mots clés du Thésaurus ou mots clés pour les metadatas du pdf
\Kone{Conception, Optimisation} % Nature de l'activité
\Ktwo{Services marchands} % Branche d'activité économique
\Kthree{Fiabilité, Informatique} % Domaines technologiques
\Kfourth{Gestionnaires de réseaux (logiciels)} % Application Physique directe
%%%%%%%%%%%%%%%%%%%%%%%%%%%%%%%%%%%%%%%%%%%%%%%%%%%%%%%%%%%%

\Semestre{Printemps 2026}
%\UE{LT01} %Nom de l'UE OU nom complet de la branche si en mode rapport de stage !
\UE{Réseaux et Télécommunications}

% Le titre de votre rapport OU le résumé de votre stage si en mode stage
% \newcommand{\titletext}{Un rapport en \LaTeX \\ écrit avec amour}

\newcommand{\titletext}{

Durant mon stage chez Winamax, leader français du poker et des paris sportifs en ligne, j'ai joué un
rôle actif dans la refonte du système d'alerting de la plateforme, chargé de prévenir les équipes
techniques lorsqu'un problème survient. Mon travail s'est inscrit dans une démarche visant à réduire le
bruit d'un système qui sollicitait inutilement les personnes d'astreinte, faute d'un manque d'alertes pertinentes et
d'une stratégie de routage fondée sur des niveaux de priorité. Après une analyse approfondie de
l'existant, j'ai conçu une solution corrigeant ces défauts. Mes missions ont notamment consisté à
mettre en place des alertes pertinentes centrées sur l'impact utilisateur, à les hiérarchiser selon leur
criticité, à enrichir leur contenu pour accélérer le diagnostic et à centraliser leur déclaration. Ce stage m'a offert l'opportunité de mettre en pratique mes
connaissances techniques tout en développant mes compétences en gestion de projet et en conception
de systèmes.
}



%% En mode Année - Mois - Jour
%\date{\today} % Pour la date de compilation
\date{\today}

\author{{\sc GONCALVES} Raphaël
% \and
% {\sc Nom} Prénom
% \break
% {\sc Nom} Prénom
% \and
% {\sc Nom} Prénom
% \break
% {\sc Nom} Prénom
}

% Texte affiché sur le carré bleu
\newcommand{\jobposition}{Ingénieur R\&T}
%\title{Vous êtes probablement assez bon pour travailler dans cette entreprise pour laquelle vous pensez ne pas être assez bon.}

\title{Amélioration du système d'alerting de Winamax}

%%%%%%% Pour les métadonnées
% Fonctionne, seulement il y a un warning à cause des "\\" dans le title
% A tester ! Pour voir les infos du pdf on fait pdfinfo sous linux
% https://ctan.gutenberg.eu.org/macros/latex/contrib/hyperref/doc/hyperref-doc.pdf#subsection.3.7
\hypersetup{
    pdfauthor = {\theauthor}, % Auteur : pdfauthor = {\theauthor}
    pdftitle = {\thetitle}, % Titre : pdftitle = {Titre du document}
    pdfkeywords = {\theKone, \theKtwo, \theKthree, \theKfourth}
}

%%%%%%%
% n'oubliez pas de changer le language principal dans rUTT.cls
% en options du package babel !

\begin{document}
    %%%% - Choix de la page de garde

    %\frontpagereports % Pour le modèle rapports de TDs / TPs / Projets
    \frontpageSTB % Pour le modèle rapports de Stages des Branches / Master
    %\frontpageSTC % Pour le modèle rapports de Stages des TC

    {
        % \selectlanguage{english}
        \myILB
    }

    % page blanche après page de garde pour impression recto verso

    \pagestyle{UTT} % Pour que notre document utilise ce style
    % Ici on organise nos parties
    \justifying % on justifie notre texte via ragged2e


    \pagenumbering{gobble} % on n'affiche pas les numéros de page
    \import{latex-files/}{remerciements.tex} % Toujours avant le sommaire !

    \clearpage

    %%%%%%%%%%%%%%%%%%%%%%
    %% Sommaire
    %%%%%%%%%%%%%%%%%%%%%%

    {
        \setcounter{tocdepth}{1}
        \renewcommand{\contentsname}{Sommaire}
        \tableofcontents
    }

    \clearpage

    % \setglossarystyle{listhypergroup}
    % \printglossary[type=\acronymtype, title={Liste des acronymes}, toctitle={Liste des acronymes}]
    % \printglossary[title={Glossaire}, toctitle={Glossaire}]

    \clearpage

    \pagenumbering{arabic}

    \import{latex-files/}{introduction.tex}

    \clearpage

    \import{latex-files/}{presentation_entreprise.tex}
    \import{latex-files/}{presentation_stage.tex}
    \import{latex-files/}{analyse_systeme_actuel.tex}
    \import{latex-files/}{conception_nouvelle_solution.tex}
    \import{latex-files/}{mise_en_oeuvre_technique.tex}
    \import{latex-files/}{experience_transverse.tex}
    \import{latex-files/}{bilan_et_perspectives.tex}
    

    \clearpage

    \import{latex-files/}{conclusion.tex}
    \label{LastPage}

    \clearpage

    % Bibliographie !
    {
        \pagenumbering{gobble} % On n'affiche pas les numéros de page
        \phantomsection % hyperlinks will target the correct page
        \markboth{Bibliographie}{}
        \raggedright % pour éviter certaines erreurs rares d'affichage
        \sloppy
        \nocite{*} % pour faire apparaître tout du fichier bib
        \printbibliography[title={Bibliographie},heading=bibintoc]
    }

    \clearpage

    %\tripleS
    \pagenumbering{Roman} % On numérote en romain pour les annexes

    % Annexes !
    \import{latex-files/annexes/}{annexes.tex}

    \clearpage
    % Toujours avoir la table des matières en dernier ! Ici figure tout, même les annexes
    % Commenter/supprimer pour enlever la table des matières
    \myfinaltoc
    % On laisse une page blanche à la fin pour l'impression, c'est plus joli
    \clearpage
    \myemptypage

\end{document}
